\chapter{A covered fallback on the same physical records}
\label{ch:og-combination}

The local-geometry certificate improves the strong-drive benchmark, but its
finite-step penalty removes all positive certificates in the weak-drive case.
The endpoint certificate retains sensitivity there. These complementary
responses suggest a practical composite inference rule. Its justification is
the established union bound; this composition introduces no new statistical
principle and requires no additional acquisition.

\section{Two certificates, one failure budget}
Let $L_G$ and $L_E$ be geometry and endpoint lower bounds computed on the
same fixed independent stationary triples, with component statistical failure
probabilities $\alpha_G$ and $\alpha_E$. Assume the two physical/calibration
contracts hold jointly. Then
\begin{equation}
 \mathbb P\{\max(L_G,L_E)>\sigma\}
 \le \mathbb P\{L_G>\sigma\}+\mathbb P\{L_E>\sigma\}
 \le \alpha_G+\alpha_E.
 \label{eq:og-combined-union}
\end{equation}
The scores may be arbitrarily dependent. The failure allocation and the two
component rules must be fixed for prospective use, or supplied by a separately
justified simultaneous procedure. Here $\alpha_G=\alpha_E=1/40$, so the covered
maximum has statistical failure probability at most $1/20$. A shared pre-test
calibration failure probability is counted once when one joint event validates
both contracts; independent calibration failures can instead be conservatively
summed before adding them. A same-test-data calibration event does not inherit
this guarantee automatically.

The geometry component uses the corrected score of
Chapter~\ref{ch:og-finite}, including the calibrated reading error and score
rounding. The endpoint component uses
\begin{equation}
 Z_s=-(sY_0-Y_{-h})(sY_0-Y_h),\qquad 0\le s\le1,
 \qquad Y_t\in[-1,1].
\end{equation}
Stationarity gives
$\mathbb E Z_s=-s^2 C(0)+2sC(h)-C(2h)$, with
$Z_s\in[-(1+s)^2,1]$. The same empirical-Bernstein construction, with its
uniform score-error correction, supplies a covered lower bound $d$ on the
positive endpoint defect. The seventh-edition physical bridge then yields
\begin{equation}
 L_E=\frac{\Phi(d)}h,\quad 0\le d<1,\qquad
 \Phi(u)=2\sqrt u\operatorname{atanh}\sqrt u,
\end{equation}
with the implementation retaining the existing rational outward enclosure and
its finite-entropy incompatibility handling. The constant follows from the
observable range $R=2$. The implementation requires the common normalized
readout contract; it does not silently use the same constants for a wider
sensor range.

\section{Retrospective reanalysis of the frozen bank}
The equal allocation was introduced after inspecting the primary benchmark.
The following is therefore a retrospective evaluation of a fixed composite
rule, not a new frozen confirmatory experiment. No timing, physical parameter,
record, generator cap or original result was changed. The newly computed
bounds use the exact stored sums and sums of squares from all $1{,}200$
datasets. The original individually valid $\alpha=0.05$ bounds remain in the
primary report; maximizing those original bounds would not by itself retain
a $0.05$ failure guarantee.

\begin{center}
\begin{tabular}{lrrr}
Scenario & Geometry at $0.025$ & Endpoint at $0.025$ & Covered maximum\\
\hline
Strong drive, $\sigma=9$ & $400/400$ & $400/400$ & $400/400$\\
Weak drive, $\sigma=1$ & $0/400$ & $383/400$ & $383/400$\\
Equilibrium, $\sigma=0$ & $0/400$ & $0/400$ & $0/400$\\
\end{tabular}
\end{center}
The geometry component wins every strong-drive dataset; its mean certified
lower bound is $2.887645$, compared with $2.982329$ at its original individual
$0.05$ allocation. The endpoint component supplies every positive weak-drive
result. Its mean bound is $0.00880994$, compared with $0.01196295$ at its
original allocation; detections fall from $396$ to $383$ out of $400$. This is
the observed cost of simultaneous coverage. No reported lower bound exceeds
the known generating entropy in these simulations, which is a numerical
check rather than a proof of coverage beyond the theorem's assumptions.

Every dataset still uses $200{,}000$ independent stationary triples,
$600{,}000$ endpoint readings and $16{,}000$ units of modeled evolution time.
Both scores reuse these same records, so those resources are counted once.
The physical cost of preparation, geometry and clock calibration, exposure,
resetting and energy remains additional. The combined result preserves the
strong-drive benefit while recovering much of the weak-drive endpoint
sensitivity. It does not dominate an individually allocated endpoint test at
every operating point, and it does not justify optimizing the allocation on
the same data without further simultaneous coverage.
